\section{Dashboard Prototyp}

\begin{figure}[h]
	\centering
	\includegraphics[width = 1.0\textwidth]{dashboard.png}
	\caption{Skizze des Dashboards}
	\label{fig:dashboard}
\end{figure}

\subsection{Überschrift}

In der Überschrift soll der User mit seinem Vornamen begrüßt und die Studienzeit in Tagen angezeigt werden, damit der/die User/User*in sieht wie lange das Studium gesamt schon läuft. Der Name soll das Dashboard etwas persönlicher gestalten.

\subsection{ECTS Counter}

Durch den ECTS Counter soll eine übersicht des Fortschritts in einer ECTS Zahl dargestellt werden.

\subsection{ECTS pro Monat}

Mit der Grafik für die ECTS pro Monat soll das Ziel mindestens 5 ECTS pro Monat zu erzielen überwacht werden. Mit dieser Grafik soll es auf einen Blick gelingen wie viele ECTS im dazugehörigen Monat erreicht wurden.

\subsection{Zeit zu ECTS}

Das Ziel das Studium in mindestens 3 Jahren abzuschließen erfordert eine Anzeige für die Zeit und eine für den Fortschritt. Diese beiden müssen gegenübergestellt werden, damit das Ziel überwacht werden kann. Beide Anzeigen zeigen die jeweiligen Werte in Prozent an, damit beides verglichen werden kann.

\subsection{Hinzufügen Button}

Durch den Button rechts unten soll es möglich sein erbrachte Prüfungsleistungen hinzuzufügen. Dabei soll in weiterer Folge ein Fenster aufgehen in dem der User die Daten für die abgeleistete Prüfungsleistung eintragen kann.